\chapter{Introduction}

\section{Mission Overview}

Lakitu is a 1U CubeSat-format payload flown aboard a stratospheric balloon,
developed as a Spacecraft Computer, Software and Data (SCSD) course project at
the Technical University of Munich. The platform executes an autonomous
4--6~hour flight managed by a six-phase finite state machine: \emph{Standby},
\emph{Launch}, \emph{Ascent}, \emph{Cruise}, \emph{Descent}, and
\emph{Landing}. Its primary payload is a downward-facing camera coupled to an
on-board machine-learning model that performs visual cloud-cover estimation in
flight. Every captured frame is archived to an SD card and a compact
cloud-fraction record is downlinked over LoRa alongside housekeeping
telemetry.

\section{Scope of this Report}

This report documents the spacecraft from a Computer, Software and Data
perspective, organised by subsystem: Command and Data Handling (CDH),
Telemetry, Tracking and Command (TT\&C), the On-Board Computer (OBC),
Electrical Power System (EPS), Flight Software (FSW) and its finite state
machine, the on-board Machine-Learning Payload, and Structures. Each chapter
covers that subsystem's responsibilities, interfaces to the rest of the
spacecraft, and key design decisions as implemented in the
\texttt{scsd-lakitu-cubesat} repository.

\section{Team and Responsibilities}

\textcolor{red}{Everyone: please check the table below is correct for your
own contributions and flag anything that needs correcting.}

\begin{table}[H]
\centering
\begin{tabular}{lll}
\toprule
\textbf{Name} & \textbf{Led} & \textbf{Helped with} \\
\midrule
Frederik Alexander Simon & CDH & -- \\
Keval Satish Mohite & Project Management, FSW & -- \\
Andrei Arteev & FSW & EPS \\
Konstantin Hesselmann & EPS & OBC, Structures \\
Alberto S\'anchez Calvo & TT\&C & -- \\
Alisher Orynbek & OBC & -- \\
Angela Ghurburrun & Structures & -- \\
\bottomrule
\end{tabular}
\caption{Team leads and supporting contributions by subsystem}
\end{table}
